\newpage
\section{Particle distance measurement}
\label{sec:cryo_fib_sem}

The inter-particle distance in fresh cement paste is crucial as it directly governs the rheology of the paste and the formation of the initial microstructure.
Tracking changes in this distance allows to quantify the effectiveness of \ac{PUS} in dispersing the cement particles.

\subsection{Sample preparation and image evaluation}

To determine the particle distance in fresh (\qty{20}{\minute}) cement pastes using a \ac{SEM}, the paste has to be frozen.
The specimen has to be flash frozen at high pressures (\qty{0.2}{\giga\pascal}) to \qty{-150}{\celsius} to avoid artefacts due to the formation of ice crystals.
Therefore, the high-pressure frozen samples of cement pastes were prepared according to \zcref{sec:cryofibnt}.
The frozen sample was coated with \ce{Pt} to prevent charging effects during the \ac{SEM} imaging.
Similar to the \ac{FIBnt} method (\zcref{sec:lit_fib}), a trench through the whole specimen thickness was cut to obtain a fresh vertical cross-section.
Due to the low temperature, the \ce{Pt} layer thickness exceeded the intended thickness of \qty{0.5}{\micro\meter} significantly (see \zcref{fig:cryo_site}).

\begin{figure}[H]
    \centering
	\includegraphics[width=\textwidth]{particle-dist/cryo_site.pdf}

    \caption{
		Overview over the \ac{FIBnt} site of the pressure-frozen \ac{C3S} paste. A: \ce{Pt} deposition, B: \ce{C3S} paste, C: brass sample holder.
	}
    \label{fig:cryo_site}
\end{figure}

Also, the redepositioning of the removed material during the cutting process was much more pronounced compared to the \ac{FIBnt} at room temperature due to the low temperature.
\ac{SEM}-\ac{BSE} images as shown in \zcref{fig:cryo-fib} were taken from the cross-section.

\begin{figure}[htb!]
    \centering
	% area: \qty{10683.3}{\micro\meter\squared}}
	\begin{subfigure}{0.49\textwidth}
		\includegraphics[width=\textwidth]{particle-dist/raw.pdf}
	\end{subfigure}
	\hfill
	\begin{subfigure}{0.49\textwidth}
		\includegraphics[width=\textwidth]{particle-dist/segmented.png}
	\end{subfigure}

    \caption{
		\ac{FIB}-section of fresh \ac{C3S}-paste, \ac{wb}\,=\,\num{0.5}, \qty{20}{\minute}.
		Denoised \ac{BSE}-image (left-hand side) and segmented image (right-hand side, {\color{red}red}: \ac{C3S}; {\color{blue}blue}: \ce{Pt}; black: ice).
	}
    \label{fig:cryo-fib}
\end{figure}

A comparison of the \ac{PSD} measured using laser diffraction and the image based analysis can be found in \zcref{fig:psd_cryo_C3S_lt20mum}.
As expected from the small and therefore not representative image area and sectioning effects (see \zcref{subsec:sectioning-effect}), the \ac{PSD} of the \ac{FIB}-\ac{SEM} image is shifted to smaller particle sizes.
To evaluate the distance between the particles within \zcref{fig:cryo-fib}, the image evaluation method used in \zcref{chap:hydrate-rim} was adapted.
Instead of measuring the distance between the particle surface and the hydrate layer surface, the distance to the next particles was measured.
Similar to the \ac{HLT} measurement, the centre of gravity of the particle was used.
The image of the centred particle was then converted to polar coordinates and the distance to the next particle was determined.

%\begin{wrapfigure}{R}{0.6\linewidth}
\begin{figure}[htb]
    %\vspace{-\intextsep}
    \centering
	\includegraphics[width=.9\linewidth, trim={.8cm 0.3cm 0.5cm 0.5cm}, clip]{particle-dist/dis_histogram.pdf}

    \caption{
		Particle distance to particle diameter distribution of \qty{20}{\minute} old \ac{C3S} paste (\ac{wb}\,=\,\num{0.5}).
	}
    \label{fig:cryo-fib-results}
    %\vspace{-\intextsep}
\end{figure}
%\end{wrapfigure}

\subsection{Dispersing effects of sonication and superplasticiser}

% \Uni\WdB\Cryo-LC3-PUS\Segmentierung LC3+PUS\
Another application of the method was used to measure the distance between particles of a CEM~II-C Q-LL, which is usually called \ac{LC3} in the literature, with and without the application of \ac{PUS}.
While the measurement of the particle size distribution based on the segmented image was already done \cite{rossler.2023,rossler.2023a} (see \zcref{fig:psd_cryo_LC3}), the distance between the particles was not measured until now.
The images have a similar size (untreated: \qty{1095.9}{\micro\meter\squared}, \ac{PUS}: \qty{1049.9}{\micro\meter\squared}) and were taken from the same sample.
As the \ac{PSD} in \zcref{fig:psd_cryo_LC3} shows, the \ac{PUS} seems to disperse the particles as expected.
The images themselves, also give the impression that the \ac{PUS} has a significant effect on the particle distance due to the more homogeneous particle distribution within the area.

\begin{figure}[htb!]
    \centering

	% area: \qty{1095.9}{\micro\meter\squared}}
	\begin{subfigure}{0.49\textwidth}
		\includegraphics[width=\textwidth]{particle-dist/LC3/noPUS/raw.pdf}
        \caption{\ac{BSE}-image of untreated \ac{LC3}-paste.}
	\end{subfigure}
	\hfill
	\begin{subfigure}{0.49\textwidth}
		\includegraphics[width=\textwidth]{particle-dist/LC3/noPUS/segmented.pdf}
        \caption{Segmented image of untreated \ac{LC3}-paste.}
	\end{subfigure}

%     \caption{
% 		\ac{FIB}-sections of fresh \ac{LC3}-paste, reference (no \ac{PUS}), \ac{wb}\,=\,\num{0.5}, \qty{20}{\minute}.
% 		Denoised \ac{BSE}-image (left-hand side) and segmented image (right-hand side, white: ignored particles; yellow: particles used for the analysis; red: pore space included in measurements; black: ignored pore space).
% 	}
%     \label{fig:cryo-fib-lc3a}
% \end{figure}

% \begin{figure}[htb!]
%     \centering
%
	% area: \qty{1049.9}{\micro\meter\squared}}
	\begin{subfigure}{0.49\textwidth}
		\includegraphics[width=\textwidth]{particle-dist/LC3/PUS/raw.pdf}
        \caption{\ac{BSE}-image of sonicated \ac{LC3}-paste.}
	\end{subfigure}
	\hfill
	\begin{subfigure}{0.49\textwidth}
		\includegraphics[width=\textwidth]{particle-dist/LC3/PUS/segmented.pdf}
        \caption{Segmented image of sonicated \ac{LC3}-paste.}
	\end{subfigure}

    \caption{
		\ac{FIB}-sections of fresh \ac{LC3}-paste, untreated (top) and sonicated with \ac{PUS} (bottom), \ac{wb}\,=\,\num{0.5}, hydrated for \qty{20}{\minute}.
		Denoised \ac{BSE}-image (left-hand side) and segmented image (right-hand side, white: ignored particles; {\color{orange}yellow}: particles used for the analysis; {\color{red}red}: pore space included in measurements; black: ignored pore space).
	}
    \label{fig:cryo-fib-lc3}
\end{figure}

The images have a very low evaluable particle count (untreated: \num{633}, \ac{PUS}: \num{436}), which is due to the small area of the \ac{FIB}-\ac{SEM} images.
This also included particles that would usually be considered to be too close to the image border.
However, some particles were excluded from the analysis due to their unfavourable particle shape.
This particularly affects the \ac{PUS}-treated sample as can be seen in \zcref{fig:cryo-fib-lc3}.
The low particle count is reflected in the low data density of the histograms in \zcref{fig:cryo-fib-lc3-results}.

\begin{figure}[htb!]
    \centering

	\begin{subfigure}{0.495\textwidth}
		\includegraphics[width=\textwidth]{particle-dist/LC3/noPUS/histogram_border.pdf}
	\end{subfigure}
	\hfill
	\begin{subfigure}{0.495\textwidth}
		\includegraphics[width=\textwidth]{particle-dist/LC3/PUS/histogram_border.pdf}
	\end{subfigure}

    \caption{
		Particle distance to particle diameter distribution of \qty{20}{\minute} old \ac{LC3} paste without (left-hand side) and with (right-hand side) the application of \ac{PUS}.
        The mean distance between particles decreased from \qty{0.31(0.16)}{\micro\meter} to \qty{0.27(0.14)}{\micro\meter}, which is not statistically significant.
        The orange histograms are directly compared in \zcref{fig:cryo-fib-results_2a}.
	}
    \label{fig:cryo-fib-lc3-results}
\end{figure}

As can be seen in \zcref{fig:cryo-fib-results_2a}, the broader particle distance distribution in the untreated specimen is particularly evident in the data.
This becomes especially clear by the higher amount of larger particle distances above \qty{1}{\micro\meter}.
This fits to the subjective appearance of the particle distribution.

\begin{figure}[H]
    \centering

	\begin{subfigure}[t]{0.495\textwidth}
		\begin{tikzpicture}
			\begin{axis}[
				width  = 1.1\textwidth,
				height = .3\textheight,
				ybar,
				bar width=.0193, % pixel width of the raw image
				xmin=0, xmax=3,
				ymin=0, ymax=1,
				legend style={
					at={(.99,.98)},
					anchor=north east,
					style={column sep=4pt}
				},
				legend columns = 1,
				ymajorticks=false,
				xlabel={particle distance in \unit{\micro\meter}},
				ylabel={normalised intensity},
				grid=both,
				grid style={line width=.1pt, draw=gray!10},
				legend cell align={left},
			]

				\addplot[fill=blue, draw=none, mark=none, opacity=0.5, draw opacity=0] table [x=X_center,y=S, col sep=comma] {figures/particle-dist/LC3/LC3_border.csv};
				\addlegendentry{ \ac{LC3} };

				\addplot[fill=red, draw=none, mark=none, opacity=0.5, draw opacity=0] table [x=X_center,y=S, col sep=comma] {figures/particle-dist/LC3/LC3+PUS_border.csv};
				\addlegendentry{ \ac{LC3}+\ac{PUS}};

				\draw[->, thick, black] (axis cs:1.2,0.4) -- ++(-0.1,-0.1)
				node[anchor=south west, font=\footnotesize, yshift=10pt] {\begin{tabular}{@{}l@{}}larger distances \\ are more prevalent \\ in the untreated specimen
				\end{tabular}};
			\end{axis}
		\end{tikzpicture}
		\caption{
            \qty{20}{\minute} old \ac{LC3} paste without ({\color{blue}blue}) and with ({\color{red}red}) the application of \ac{PUS}.
        } %  The mean distance between particles decreased from \qty{0.33(0.17)}{\micro\meter} and \qty{0.24(0.15)}{\micro\meter}.
		\label{fig:cryo-fib-results_2a}
	\end{subfigure}
	\hfill
	\begin{subfigure}[t]{0.495\textwidth}
		\begin{tikzpicture}
			\begin{axis}[
				width  = 1.1\textwidth,
				height = .3\textheight,
				ybar,
				bar width=.030304, % pixel width of the raw image
				xmin=0, xmax=3,
				ymin=0, ymax=1,
				legend style={
					at={(.99,.98)},
					anchor=north east,
					style={column sep=4pt}
				},
				legend columns = 1,
				ymajorticks=false,
				xlabel={particle distance in \unit{\micro\meter}},
				ylabel={normalised intensity},
				grid=both,
				grid style={line width=.1pt, draw=gray!10},
				yticklabel pos=right,
				legend cell align={left},
			]

				\addplot[fill=blue, draw=none, mark=none, opacity=0.5, draw opacity=0] table [x=X_center,y=S, col sep=comma] {figures/particle-dist/zing.2008a/f5_histogram.csv};
				\addlegendentry{ \acs*{OPC}, no \acs*{PCE}};

				\addplot[fill=red, draw=none, mark=none, opacity=0.5, draw opacity=0] table [x=X_center,y=S, col sep=comma] {figures/particle-dist/zing.2008a/f7_histogram.csv};
				\addlegendentry{ \acs*{OPC}, \acs*{PCE}};

				\draw[->, thick, black] (axis cs:1.2,0.48) -- ++(-0.1,-0.1)
				node[anchor=south west, font=\footnotesize, yshift=10pt] {\begin{tabular}{@{}l@{}}larger distances \\ are more prevalent \\ in the untreated specimen\end{tabular}};
			\end{axis}
		\end{tikzpicture}
		\caption{
            Fresh \ac{OPC} paste without ({\color{blue}blue}) and with ({\color{red}red}) \acs*{PCE}.
		    Source images from \cite{zingg.2008a}.
        }
		\label{fig:cryo-fib-results_2b}
	\end{subfigure}

    \caption{
		Particle distance distributions of \acs*{LC3} and \acs*{OPC} paste comparing the influence of \acs*{PUS} and \acs*{PCE} respectively.
		Both methods shift the particle distance distribution to smaller distances, which means they have a dispersing effect.
    }
    \label{fig:cryo-fib-results_2}
\end{figure}

These data illustrate that \ac{PUS} is able to disperse the particles of the \ac{LC3} more homogeneously.
However, a side effect of the still agglomerated particles is their reduced number.
Furthermore, the image data acquired by cryo-\ac{FIBnt} are of limited significance due to the small imaged area and therefore limited representativeness.

Due to the expensive experiments, two figures published by \textcite{zingg.2008a} (Figures~\num{5} and \num{7}) were also used to test the method.
They used cryo-\ac{FIB}-\ac{SEM} to image the influence of \ac{PCE} on the fine fraction of \ac{OPC} in fresh cement pastes (\ac{wb}\,=\,\num{0.5}).
The resulting images show a distinct difference in the particle distribution.
The images were processed similar to the \ac{LC3} specimens.

The particle distance distributions for the \ac{OPC}+\ac{PCE} experiments from \textcite{zingg.2008a} are shown in \zcref{fig:cryo-fib-results_2b} and show a more significant difference in the particle distribution.
The untreated \ac{OPC} shows more large particle distances than the specimen mixed with \ac{PCE}.
The segmentation results can be found in \zcref{fig:cryo-fib-zingg-results} and the particle distance to particle diameter distributions in \zcref{fig:cryo-fib-zingg-results2}.

The results of both test series show that \ac{PUS} and \ac{PCE} have, as expected, a dispersing effect on the particles in fresh cement pastes, which can be visualised and evaluated objectively by means of cryo-\ac{FIB}-\ac{SEM}. % SRC

% \FloatBarrier
% \section{Conclusion}

% The image evaluation method developed for \ac{HLT} measurements in \zcref{chap:hydrate-rim} was successfully adapted to quantify inter-particle distances in fresh cement paste from cryo-\ac{FIB}-\ac{SEM} images.
% Applied to \ac{C3S} paste, \ac{LC3} with and without the application of \ac{PUS}, and to published \ac{OPC}+\ac{PCE} data, the method confirmed that both \ac{PUS} and \ac{PCE} shift the particle distance distribution towards smaller values and thus have a dispersing effect.
% The broader distribution of particle distances in untreated specimens, compared to sonicated or superplasticised ones, is reflected in the histograms and supports the qualitative impression from the images.
% The main limitation remains the small imaged area of cryo-\ac{FIB} sections, which restricts statistical representativeness and particle counts; the approach is nevertheless suitable to objectify dispersion effects in fresh paste and to complement rheological or other bulk measurements.

% image preparation for images from \cite{zingg.2008a}:
% -> manual background removal (shadowing)
% -> denoising (despecle) using \textsc{ImageJ}
% -> segmentation using \textsc{ImageJ} (thresholding+watershed)
% -> application of the particle distance measurement method from \zcref{chap:hydrate-rim} (see \zcref{sec:cryo_fib_sem}) to the segmented image